\documentclass[10pt,a4paper]{article}
\usepackage{amsmath,amssymb,amsthm,graphicx}
\usepackage{amsmath,amssymb,amsthm}
\newcommand{\FF}{\mathbb F}
\newcommand{\PP}{\mathbb P}
\newcommand{\Aspace}{\mathbb A}
\newcommand{\GG}{\mathrm G}
\newcommand{\supp}{\operatorname{supp}}
\newenvironment{result}[1]{\par\medskip\noindent\textbf{#1.}\itshape}{\par\medskip}
\title{Chow-variety dimension lower bounds on the size of variety-evasive subspace families}
\author{Zeyu Guo}
\date{doi:10.1007/s00037-024-00256-1; CC BY 4.0}
\begin{document}
\maketitle
\noindent\textit{Editorial note.} Faithful editable transcription of the actual complete CC BY2024 journal main lower-bound proof and necessary original context. Original theorem/page references retained; classical geometry facts are standard prerequisites with the original author-provided statements/citations. No proof substituted or mathematical correction made. Native TeX retrieval has remote connection timeouts; actual journalPDF/page already retained, so this record is based directly on the licensed journal text. Independent lower-bound mechanism counted once, without splitting degree cases or classical Chow-dimension support. All production/extraction/validation solely on Yupeng.\par\medskip
\noindent\textit{Allowed standard prerequisites (editorial).}\begin{itemize}
\item Projective algebraic varieties, Grassmannians/Plucker coordinates, Chow embeddings and the standard dimension formula for Chow varieties. Original definitions and citedTheorems3.5/3.7/4.1/4.4 retained, with author-provided explanatory argument. Classical results and projective hyperplane dimension fact cited rather than recursively reproved.
\end{itemize}\medskip
\section*{Original ambient variety and evasion definitions (journal pp.3--6)}
Let $\FF$ be an algebraically closed field. An affine $n$-space $\Aspace^n$, as a set, is simply defined to be the vector space $\FF^n$. We also need the notion of a projective $n$-space, denoted by $\PP^n$, which is (intuitively) the set of lines passing through the origin $0$ of $\Aspace^{n+1}$. Formally, it is defined to be the quotient set $(\Aspace^{n+1}\setminus\{0\})/\sim$, where $\sim$ is the equivalence relation defined by scaling, i.e., $u\sim v$ if $u=cv$ for some nonzero scalar $c\in\FF$.
An (affine) subvariety $\mathcal V\subseteq\Aspace^n$ is the set of common zeros of a set of $n$-variate polynomials over $\FF$. Similarly, a (projective) subvariety $\mathcal V\subseteq\PP^n$ is the set of common zeros of a set of homogeneous $(n+1)$-variate polynomials over $\FF$, where we represent each element of $\PP^n$ as an $(n+1)$-tuple in $\Aspace^{n+1}$. In this paper, a variety refers to a subvariety of a projective or affine space, and is said to be irreducible if it cannot be written as a union of finitely many proper subvarieties. (Varieties in this paper are not necessarily irreducible. They are often called algebraic sets in literature.)
The dimension of a variety $\mathcal V$, denoted by $\dim(\mathcal V)$, is intuitively the ``degree of freedom'' of picking a point in the variety. See Section 2.3 for its formal definition. In particular, for a linear subspace $\mathcal V\subseteq\Aspace^n$, the dimension of $\mathcal V$ as a variety is simply its linear-algebraic dimension.
\begin{result}{Definition 1.1 (Evading)}
Let $\mathcal V_1$ and $\mathcal V_2$ be irreducible subvarieties of $\PP^n$ or $\Aspace^n$. We say $\mathcal V_1$ evades $\mathcal V_2$ if
\[\dim(\mathcal V_1\cap\mathcal V_2)\le\dim(\mathcal V_1)+\dim(\mathcal V_2)-n,\]
where the dimension of an empty set is assumed to be $-\infty$. In particular, if $\dim(\mathcal V_1)+\dim(\mathcal V_2)<n$, then $\mathcal V_1$ evades $\mathcal V_2$ iff $\mathcal V_1\cap\mathcal V_2=\emptyset$.
More generally, suppose $\mathcal V_1$ is irreducible but $\mathcal V_2$ is possibly reducible. We say $\mathcal V_1$ evades $\mathcal V_2$ if it evades every irreducible component of $\mathcal V_2$.
\end{result}
\begin{result}{Definition 1.2 (Subspace family)}
For $0\le k\le n$, a finite collection (i.e., multiset) of $k$-subspaces of $\PP^n$ is called a (projective) $k$-subspace family on $\PP^n$. Similarly, a finite collection of affine $k$-subspaces of $\Aspace^n$ is called an affine $k$-subspace family on $\Aspace^n$.
\end{result}
\begin{result}{Definition 1.3 (Variety evasive subspace family)}
Let $\mathcal F$ be a family of subvarieties of $\PP^n$ (resp. $\Aspace^n$). Let $\mathcal H$ be a $k$-subspace family on $\PP^n$ (resp. affine $k$-subspace family on $\Aspace^n$) where $0\le k\le n$. Then:
\begin{itemize}\item We say $\mathcal H$ is $\mathcal F$-evasive if for every $\mathcal V\in\mathcal F$, there exists $W\in\mathcal H$ that evades $\mathcal V$.
\item We say $\mathcal H$ is $(\mathcal F,\epsilon)$-evasive if for every $\mathcal V\in\mathcal F$, a random element $W\in\mathcal H$ evades $\mathcal V$ with probability at least $1-\epsilon$.\end{itemize}
\end{result}
\begin{result}{Definition 1.4 (degree)}
The degree of an irreducible variety $\mathcal V$ in $\PP^n$ (resp. $\Aspace^n$) is the number of intersections of $\mathcal V$ with a general projective (resp. affine) subspace of codimension $\dim(\mathcal V)$. Following Heintz (1983); Heintz \& Schnorr (1980), we define the degree of a (possibly reducible) variety to be the sum of the degrees of its irreducible components.
\end{result}
\section*{Original Pl\"ucker and Chow-form background (journal pp.27--30)}
\paragraph{Grassmannians.}
Let $n\in\mathbb N$ and $k\in\{0,1,\ldots,n-1\}$. The Grassmannian $\GG(k+1,n+1)$ is the set of all $(k+1)$-dimensional linear subspaces of $\Aspace^{n+1}$. By taking the quotient modulo scalars, it may also be identified with the set of all $k$-subspaces of $\PP^n$, which we denote by $\GG(k,n)$.
\paragraph{The Pl\"ucker embedding and Pl\"ucker coordinates.}
Consider a linear subspace $W\in\GG(k+1,n+1)$. The simplest way of representing $W$ is using a $(k+1)\times(n+1)$ matrix $A$ over $\FF$ such that $W$ equals the row space of $A$. We call such a matrix $A$ a generating matrix of $W$. For convenience, we also say $A$ is a generating matrix of $\PP W\in\GG(k,n)$.
\begin{result}{Definition 3.4 (Pl\"ucker embedding)}
Define $\phi:\GG(k+1,n+1)\to\PP^{\binom{n+1}{k+1}-1}$ by
\[\phi(W)=(\det A_{[k+1],S})_{S\in\binom{[n+1]}{k+1}}\]
where $A$ is a generating matrix of $W$.
\end{result}
The Pl\"ucker embedding embeds the Grassmannian $\GG(k+1,n+1)$ in $\PP^{\binom{n+1}{k+1}-1}$ as an irreducible projective subvariety, as stated by the following theorem. See, e.g., Fulton (1997); Harris (1992) for proofs.
\begin{result}{Theorem 3.5}
The Pl\"ucker embedding $\phi$ is a well-defined injective map whose image is an irreducible projective subvariety of $\PP^{\binom{n+1}{k+1}-1}$.
\end{result}
The homogeneous coordinates $(\det A_{[k+1],S})_{S\in\binom{[n+1]}{k+1}}$ of $\phi(W)$ are called the (primal) Pl\"ucker coordinates of $W$.
Denote by $R:=\FF[X_S:S\in\binom{[n+1]}{k+1}]$ the homogeneous coordinate ring of $\PP^{\binom{n+1}{k+1}-1}$. The irreducible projective subvariety $\phi(\GG(k+1,n+1))$ is defined by a homogeneous prime ideal of $R$, which we denoted by $I$. Then $R/I$ is the homogeneous coordinate ring of $\phi(\GG(k+1,n+1))$. The ideal $I$ contains precisely the polynomial relations that the Pl\"ucker coordinates need to satisfy. It is also known that $I$ is generated by certain quadratic forms, known as the Pl\"ucker relations. See Fulton (1997); Harris (1992) for details.
\paragraph{Dual Pl\"ucker coordinates.}
Alternatively, we could represent a linear subspace $W\in\GG(k+1,n+1)$ by an $(n-k)\times(n+1)$ matrix $B$ over $\FF$ whose rows specify the linear equations defining $W$. We call such a matrix $B$ a parity check matrix of $W$. For convenience, we also say $B$ is a parity check matrix of $\PP W\in\GG(k,n)$.
The entries of $B$ are called the dual Stiefel coordinates of $W$. This gives another embedding $\phi^\vee:\GG(k+1,n+1)\to\PP^{\binom{n+1}{n-k}-1}=\PP^{\binom{n+1}{k+1}-1}$, defined by
\[\phi^\vee(W)=(\det B_{[n-k],S})_{S\in\binom{[n+1]}{n-k}}.\]
The homogeneous coordinates $(\det B_{[n-k],S})_{S\in\binom{[n+1]}{n-k}}$ of $\phi^\vee(W)$ are called the dual Pl\"ucker coordinates of $W$.
Recall that we denote by $\GG(k,n)$ the set of all $k$-subspaces of $\PP^n$. By identifying $\GG(k+1,n+1)$ with $\GG(k,n)$ via $W\mapsto\PP W$, we regard $\phi$ and $\phi^\vee$ as maps from $\GG(k,n)$ to $\PP^{\binom{n+1}{k+1}-1}$.
\begin{result}{Definition 3.6 (Gelfand et al. 1994)}
For an irreducible subvariety $\mathcal V\subseteq\PP^n$ of dimension $n-k-1$, define the associated hypersurfaces $Z_{\mathcal V}$ of $\mathcal V$ as the set of $k$-subspaces intersecting $\mathcal V$, i.e.,
\[Z_{\mathcal V}:=\{W\in\GG(k,n):\mathcal V\cap W\ne\emptyset\}.\]
\end{result}
The term ``associated hypersurface'' is justified by the following theorem.
\begin{result}{Theorem 3.7}
Let $\mathcal V\subseteq\PP^n$ be an irreducible projective subvariety of dimension $n-k-1$ and degree $d\in\mathbb N^+$. Then there exists a nonzero homogeneous polynomial $P_{\mathcal V}\in R=\FF[X_S:S\in\binom{[n+1]}{k+1}]$ of degree $d$ such that $\phi(Z_{\mathcal V})$ is defined by $P_{\mathcal V}$ as a subvariety of $\phi(\GG(k,n))$. That is,
\[\phi(Z_{\mathcal V})=\phi(\GG(k,n))\cap V(P_{\mathcal V}).\]
Moreover, $R_{\mathcal V}:=P_{\mathcal V}+I\in(R/I)_d$ is uniquely determined by $\mathcal V$ up to scalars.
\end{result}
Theorem 3.7 is explicitly stated as Dalbec \& Sturmfels (1995, Theorem 1.1 and Corollary 2.1). A proof can be found in Gelfand et al. (1994, Section 3.2). We briefly explain how to find a polynomial $P_{\mathcal V}$ satisfying Theorem 3.7: Firstly, it can be shown using the trick of dimension counting via incidence varieties that $\phi(Z_{\mathcal V})$ is an irreducible projective subvariety of the Grassmannian $\phi(\GG(k,n))$ of codimension one (Gelfand et al. 1994, Section 3.2, Proposition 2.2). Secondly, the homogeneous coordinate ring $R/I$ of the Grassmannian is known to be a unique factorization domain (Fulton 1997, Chapter 9). These two facts imply that the homogeneous ideal of $R/I$ defining $\phi(Z_{\mathcal V})$ is a principal ideal. Choose $R_{\mathcal V}$ to be a generator of this principal ideal, which is unique up to scalars. Then lift $R_{\mathcal V}\in R/I$ to $P_{\mathcal V}\in R$.
\begin{result}{Definition 3.8 (Chow form)}
Let $\mathcal V\subseteq\PP^n$ be an irreducible subvariety of dimension $n-k-1$ and degree $d\in\mathbb N^+$. Define the Chow form of $\mathcal V$ in Pl\"ucker coordinates, or simply the Chow form of $\mathcal V$, to be $R_{\mathcal V}\in(R/I)_d$ as in Theorem 3.7.
More generally, for an equidimensional subvariety $\mathcal V=\bigcup_{i=1}^s\mathcal V_i\subseteq\PP^n$ of dimension $n-k-1$ and degree $d$, where $\mathcal V_1,\ldots,\mathcal V_s$ are the irreducible components of $\mathcal V$, the Chow form of $\mathcal V$ is $R_{\mathcal V}:=\prod_{i=1}^sR_{\mathcal V_i}\in(R/I)_d$. It is uniquely determined by $\mathcal V$ up to scalars.
\end{result}
As a $k$-subspace intersects $\mathcal V=\bigcup_{i=1}^s\mathcal V_i$ iff it intersects some $\mathcal V_i$, we see from Theorem 3.7 that the Chow form $R_{\mathcal V}$ of an equidimensional projective subvariety $\mathcal V$ of dimension $n-k-1$ vanishes precisely at the set of $k$-subspaces that intersect $\mathcal V$.

\begin{result}{Theorem 1.10 (journal pp.8,43--44)}
Let $n,d\in\mathbb N^+$ and $k\in\{0,1,\ldots,n-1\}$. Let $\mathcal F$ be the family of equidimensional projective subvarieties of $\PP^n$ of dimension $n-k-1$ and degree at most $d$. Suppose $\mathcal H$ is an $\mathcal F$-evasive $k$-subspace family on $\PP^n$. Then
\[
|\mathcal H|\ge\begin{cases}
(n-k)(k+1)+1&\text{if }d=1,\\[2pt]
\max\left\{d(n-k)(k+1)+1,\binom{d+n-k}{d}+(n-k+1)k\right\}&\text{if }d>1.
\end{cases}
\]
In particular, $|\mathcal H|$ is superpolynomial in $n$ when $n-k=n^{\Omega(1)}$ and $d=\omega(1)$.
\end{result}

\section*{Original Section 4: Lower Bound (journal pp.41--44)}
We prove Theorem 1.10 in this section. The main tool is the notion of Chow varieties, which parameterize projective subvarieties. More precisely, they parametrize a generalization of projective subvarieties, called (effective) algebraic cycles on a projective space.
\paragraph{Algebraic cycles.}
An algebraic $r$-cycle (or simply $r$-cycle) on $\PP^n$ is a formal linear combination $\mathcal D=\sum c_i\mathcal V_i$ of finitely many irreducible subvarieties $\mathcal V_i\subseteq\PP^n$ of dimension $r$, where the coefficients $c_i$ are integers. The degree of $\mathcal D$ is $\deg(\mathcal D):=\sum c_i\deg(\mathcal V_i)$. The support of $\mathcal D$ is $\supp(\mathcal D):=\bigcup_{c_i\ne0}\mathcal V_i$. An $r$-cycle is effective if all its coefficients are nonnegative. Denote by $C(r,d,n)$ the set of all effective $r$-cycles of degree $d$ on $\PP^n$.
\paragraph{Chow varieties.}
Let $k\in\{0,1,\ldots,n-1\}$ and $r=n-k-1$. The definition of Chow forms naturally extends to effective $r$-cycles. Namely, for an effective $r$-cycle $\mathcal D=\sum_{i=1}^r c_i\mathcal V_i$ of degree $d$ on $\PP^n$, define the Chow form of $\mathcal D$ to be $R_{\mathcal D}:=\prod_{i=1}^r R_{\mathcal V_i}^{c_i}$.
Note that $R_{\mathcal D}$ is a vector in $(R/I)_d$ and is uniquely determined by $\mathcal D$ up to scalars. Write $[R_{\mathcal D}]$ for the point in $\PP(R/I)_d$ represented by $R_{\mathcal D}$. Then we have map $\psi:C(r,d,n)\to\PP(R/I)_d$, given by
\[\psi:\mathcal D\mapsto[R_{\mathcal D}],\]
called the Chow embedding of $C(r,d,n)$. Indeed, it embeds $C(r,d,n)$ in $\PP(R/I)_d$ as a projective subvariety, as stated by the following theorem.
\begin{result}{Theorem 4.1 (Chow \& van der Waerden 1937)}
The map $\psi$ is injective and its image is Zariski-closed.
\end{result}
A proof can also be found in Gelfand et al. (1994, Chapter 4). We identify $C(r,d,n)$ with its image under $\psi$ and view it as a projective variety. This variety is called the Chow variety of effective $r$-cycles of degree $d$ on $\PP^n$.
\begin{result}{Example 4.2}
Let $V$ be the subspace of homogeneous polynomials in $\FF[X_1,\ldots,X_{n+1}]$ of degree $d$. Then $C(n-1,d,n)$ is simply the projective space $\PP V$ (see Example 3.9).
\end{result}
\begin{result}{Example 4.3}
$C(r,1,n)$ is the Grassmannian $\GG(r+1,n+1)$ (or $\GG(r,n)$) embedded in $\PP^{\binom{n+1}{r+1}-1}=\PP^{\binom{n+1}{k+1}-1}$ via $\phi^\vee$ (see Example 3.10).
\end{result}
\paragraph{The dimension of Chow varieties.}
When $d=1$, the Chow variety $C(r,d,n)$ is just the Grassmannian $\GG(r+1,n+1)$ (see Example 4.3) and its dimension is well known to be $(r+1)(n-r)$ (Harris 1992). When $d>1$, the dimension of $C(r,d,n)$ was determined by Azcue (1992) and independently by Lehmann (2017). We state their result as follows.
\begin{result}{Theorem 4.4 (Azcue 1992; Lehmann 2017)}
For $d>1$ and $0\le r<n$, the dimension of $C(r,d,n)$ is
\[\max\left\{d(r+1)(n-r),\binom{d+r+1}{r+1}-1+(r+2)(n-r-1)\right\}.\]
\end{result}
This theorem was previously proved by Eisenbud \& Harris (1992) for the special case $r=1$.
\begin{result}{Remark 4.5}
To prove Theorem 1.10, we only need a lower bound for the dimension of the Chow variety, which is much easier to prove than Theorem 4.4. Indeed, it is not difficult to see that $d(r+1)(n-r)$ is the dimension of the space of unions of $d$ $r$-subspaces of $\PP^n$, and $\binom{d+r+1}{r+1}-1+(r+2)(n-r-1)$ is the dimension of the space of degree-$d$ hypersurfaces in $(r+1)$-subspaces of $\PP^n$.
\end{result}
\paragraph{Lower bound via dimension counting.}
We now restate Theorem 1.10 and prove it using a dimension counting argument.
\begin{proof}[Proof of Theorem 1.10]
Consider an arbitrary $k$-subspace $W\in\mathcal H$. We may think of each point in $\PP(R/I)_d$ as a homogeneous polynomial of degree $d$ in Pl\"ucker coordinates modulo scalars and the ideal $I$ of Pl\"ucker relations. We know Pl\"ucker coordinates always satisfy the Pl\"ucker relations. So it makes sense to talk about if a point in $\PP(R/I)_d$ vanishes at $\phi(W)$ or not, as it does not depend on the choice of the homogeneous polynomial representing this point. Note that the constraint of $p\in\PP(R/I)_d$ vanishing at $\phi(W)$ is a linear equation in the homogeneous coordinates of $p$. So the set of points in $\PP(R/I)_d$ vanishing at $\phi(W)$ is a hyperplane of $\PP(R/I)_d$, which we denote by $H_W$.
Let $r=n-k-1$. Assume $|\mathcal H|\le\dim(C(r,d,n))$. Then we have
\[C(r,d,n)\cap\bigcap_{W\in\mathcal H}H_W\ne\emptyset\]
since taking the intersection with a hyperplane reduces the dimension of a projective subvariety by at most one (Lemma 2.7 or Lemma 2.9). So there exists an effective $r$-cycle $\mathcal D=\sum_{i=1}^s c_i\mathcal V_i\in C(r,d,n)$, where $c_1,\ldots,c_s>0$, such that $\psi(\mathcal D)=[R_{\mathcal D}]$ vanishes at $\phi(W)$ for all $W\in\mathcal H$.
Let $\mathcal V=\supp(\mathcal D)=\bigcup_{i=1}^s\mathcal V_i$. Note $\mathcal V\in\mathcal F$ since $\deg(\mathcal V)\le\deg(\mathcal D)=d$. For all $W\in\mathcal H$, we know $R_{\mathcal D}=\prod_{i=1}^sR_{\mathcal V_i}^{c_i}$ vanishes at $\phi(W)$, or equivalently, $R_{\mathcal V}=\prod_{i=1}^sR_{\mathcal V_i}$ vanishes at $\phi(W)$. This implies $\mathcal V\cap W\ne\emptyset$ for all $W\in\mathcal H$. As $\mathcal V\in\mathcal F$, this contradicts our assumption about $\mathcal H$. We conclude
\[|\mathcal H|\ge\dim(C(r,d,n))+1.\]
The dimension of $C(r,d,n)$ is $(r+1)(n-r)$ when $d=1$ and is given by Theorem 4.4 when $d>1$. Plugging in $r=n-k-1$ proves the theorem.
\end{proof}

\begin{thebibliography}{99}
\bibitem{Azcue1992} Pablo Azcue. \emph{On the dimension of the Chow varieties}. Ph.D. thesis, Harvard University, 1992.
\bibitem{Chow1937} Wei-Liang Chow and B.L. van der Waerden. Zur algebraischen Geometrie. IX. \emph{Mathematische Annalen}113(1),692--704,1937.
\bibitem{Dalbec1995} John Dalbec and Bernd Sturmfels. Introduction to Chow forms. \emph{Invariant Methods in Discrete and Computational Geometry},37--58,1995.
\bibitem{Eisenbud1992} David Eisenbud and Joe Harris. The dimension of the Chow variety of curves. \emph{Compositio Mathematica}83(3),291--310,1992.
\bibitem{Fulton1997} William Fulton. \emph{Young Tableaux: With Applications to Representation Theory and Geometry},volume35. Cambridge University Press,1997.
\bibitem{Gelfand1994} Israel M. Gelfand, Mikhail M. Kapranov and Andrei V. Zelevinsky. \emph{Discriminants, Resultants and Multidimensional Determinants}. Birkh\"auser,1994.
\bibitem{Harris1992} Joe Harris. \emph{Algebraic Geometry: A First Course}. Springer-Verlag,1992.
\bibitem{Heintz1983} Joos Heintz. Definability and fast quantifier elimination in algebraically closed fields. \emph{Theoretical Computer Science}24(3),239--277,1983.
\bibitem{Heintz1980} Joos Heintz and Claus-Peter Schnorr. Testing polynomials which are easy to compute (Extended Abstract). In \emph{Proceedings of the12th annual ACM Symposium on Theory of Computing},262--272,ACM,1980.
\bibitem{Lehmann2017} Brian Lehmann. Asymptotic behavior of the dimension of the Chow variety. \emph{Advances in Mathematics}308,815--835,2017.
\bibitem{Shafarevich1994} Igor R. Shafarevich. \emph{Basic Algebraic Geometry1: Varieties in Projective Space}. Springer-Verlag,1994.
\end{thebibliography}
\end{document}
